\documentclass[11pt]{article}
\usepackage{fontspec}
\setmainfont{Times New Roman}
\setsansfont{Arial}
\setmonofont{Consolas}
\usepackage{amsmath,amssymb}
\usepackage{geometry}
\usepackage{microtype}
\geometry{a4paper,margin=3cm}
\setlength{\parindent}{1.25em}
\setlength{\parskip}{0pt}
\emergencystretch=4em
\allowdisplaybreaks
\let\A\relax
\newcommand{\A}{\mathfrak A}
\newcommand{\R}{\mathfrak R}

\begin{document}

\begin{center}
\textbf{\S\ 15. \quad Формы десятого поряда (Систем III).}
\end{center}
\[
\text{Свијанье }\left\{\begin{array}{l}
\theta^2,\\
f\cdot j \text{ с }L,\\
\theta \text{ с }H^2.
\end{array}\right.
\]

\noindent\textit{Нередуцибилне формы:}
\[
\begin{array}{c@{\quad}l@{\qquad}c@{\quad}l}
\mathrm{A)}&(\vartheta^2lx)^4,\quad
\theta_i(\vartheta^2lx)^4&
\mathrm{B)}&(aLu)^2L_j,\quad a_i(aLu)^2L_j,\\[0.8em]
\mathrm{C)}&(\theta H^2u)^3,\quad
(\theta H^2u)^4,\quad
H_\vartheta(\theta H^2u),\quad
\theta_\eta(\theta H^2u)^3.
\end{array}
\]

\noindent\textit{Редукције:}

\noindent\mbox{A)} и \mbox{B)}: остале формы разкладајут се чрез једноразово свијанье с разложимыми формами.

\noindent\mbox{C)} Остале формы разкладајут се по \S\ 13 \mbox{B)} чрез дуалистично противоположно разматрјенье.

\noindent\textit{Последки:}

$\theta^2K$ не входи в свијанье с $L$; соответно $H^2L$ не входи в свијанье с $K$. $H^3$ и $H^2L$ не входет в свијанье с $\theta$. Тым је схема свијаньа из \S\ 7 доказана.

\end{document}

